% Part: incompleteness
% Chapter: arithmetization-syntax
% Section: coding-symbols

\documentclass[../../../include/open-logic-section]{subfiles}

\begin{document}

\olfileid[pt]{inc}{art}{cod}
\olsection{Codificação de símbolos}

A linguagem básica~$\Lang L$ da lógica de primeira ordem utiliza os símbolos
\[
\lfalse \quad \lnot \quad \lor \quad \land \quad \lif \quad \lforall
\quad \lexists \quad \eq \quad ( \quad ) \quad ,
\]
juntamente com conjuntos !!{enumerable} de variáveis e !!{constant}s e
com conjuntos !!{enumerable} de !!{function}s e !!{predicate}s de
aridade arbitrária. Podemos atribuir \emph{códigos} a cada um desses
símbolos de tal maneira que cada símbolo receba um único número como
código e dois símbolos diferentes nunca recebam o mesmo número. Sabemos
que isso é possível porque o conjunto de todos os símbolos é
!!{enumerable}; portanto, existe !!a{bijection} entre ele e o conjunto
dos números naturais. Mas queremos garantir que seja possível recuperar
o símbolo (bem como alguma informação sobre ele, por exemplo, a aridade
de !!a{function}) a partir de seu código de maneira computável. Há,
naturalmente, muitas maneiras possíveis de fazer isso. Eis uma delas,
que usa funções recursivas primitivas. (Lembre-se de que
$\tuple{n_0, \dots, n_k}$ é o número que codifica a sequência de
números $n_0$, \dots, $n_k$.)

\begin{defn}
Se $s$ é um símbolo de~$\Lang L$, defina o \emph{código do
  símbolo}~$\scode s$ da seguinte maneira:
\begin{enumerate}
\item Se $s$ está entre os símbolos lógicos, $\scode s$ é dado pela
  tabela a seguir:
\[
\begin{array}{cccccccccc}
  \lfalse & \lnot & \lor & \land & \lif & \lforall \\
\tuple{0, 0} & \tuple{0, 1} & \tuple{0, 2} & \tuple{0, 3} &
\tuple{0, 4} & \tuple{0, 5} \\
\lexists & \eq & ( & ) & ,\\
\tuple{0, 6} & \tuple{0, 7} &
\tuple{0, 8} & \tuple{0, 9} & \tuple{0, 10}
\end{array}
\]
\item Se $s$ é a $i$-ésima variável $\Obj v_i$, então
  $\scode s = \tuple{1, i}$.
\item Se $s$ é a $i$-ésima !!{constant}~$\Obj c_i$, então
  $\scode s = \tuple{2, i}$.
\item Se $s$ é a $i$-ésima !!{function} $n$-ária~$\Obj f_i^n$, então
  $\scode s = \tuple{3, n, i}$.
\item Se $s$ é o $i$-ésimo !!{predicate} $n$-ário~$\Obj P_i^n$, então
  $\scode s = \tuple{4, n, i}$.
\end{enumerate}
\end{defn}

\begin{prop}
As relações a seguir são recursivas primitivas:
\begin{enumerate}
\item $\fn{Fn}(x, n)$ sse $x$ é o código de $\Obj f^n_i$ para algum~$i$,
  isto é, $x$ é o código de !!a{function} $n$-ária.
\item $\fn{Pred}(x, n)$ sse $x$ é o código de $\Obj P^n_i$ para
  algum~$i$, ou $x$ é o código de $\eq$ e $n = 2$, isto é, $x$ é o
  código de !!a{predicate} $n$-ário.
\end{enumerate}
\end{prop}

\begin{defn}
Se $s_0, \dots, s_{n-1}$ é uma sequência de símbolos, seu \emph{número
  de G\"odel} é $\tuple{\scode{s_0}, \dots, \scode{s_{n-1}}}$.
\end{defn}

\begin{explain}
Observe que \emph{códigos} e \emph{números de G\"odel} são coisas
diferentes. Por exemplo, a variável~$\Obj v_5$ tem o código~$\scode{\Obj
  v_5} = \tuple{1, 5} = 2^2\cdot 3^6$. Mas a variável~$\Obj v_5$,
considerada como termo, também é uma sequência de símbolos (de
comprimento~$1$). O \emph{número de G\"odel}~$\Gn{\Obj v_5}$ do
\emph{termo}~$\Obj v_5$ é $\tuple{\scode{\Obj v_5}} = 2^{\scode{\Obj
    v_5} + 1} = 2^{2^2\cdot 3^6 + 1}$.
\end{explain}

\begin{ex}
Lembre-se de que, se $k_0$, \dots, $k_{n-1}$ é uma sequência de
números, então o código da sequência $\tuple{k_0, \dots, k_{n-1}}$ na
codificação por potências de primos é
\[
2^{k_0+1}\cdot3^{k_1+1}\cdot \dots \cdot p_{n-1}^{k_{n-1}+1},
\]
em que $p_i$ é o $i$-ésimo primo (começando com $p_0 = 2$). Assim, por
exemplo, a fórmula $\eq[\Obj v_0][\Obj 0]$, ou, mais explicitamente,
${\eq}(\Obj v_0,\Obj c_0)$, tem o número de G\"odel
\[
\tuple{\scode{\eq},\scode{(},\scode{\Obj v_0},\scode{,}, \scode{\Obj
    c_0},\scode{)}}.
\]
Aqui, $\scode{\eq}$ é $\tuple{0,7} = 2^{0+1}\cdot
3^{7+1}$, $\scode{\Obj v_0}$ é $\tuple{1,0} = 2^{1+1}\cdot3^{0+1}$,
etc. Portanto, $\Gn{=(\Obj v_0,\Obj c_0)}$ é
\begin{multline*}
2^{\scode{=} + 1}\cdot 3^{\scode{(}+1}\cdot 5^{\scode{\Obj v_0}+1}
\cdot 7^{\scode{,} + 1} \cdot 11^{\scode{\Obj c_0}+1} \cdot
13^{\scode{)}+1} = \\
2^{2^1\cdot 3^8 + 1}\cdot 3^{2^1\cdot 3^9+1}\cdot 5^{2^2\cdot 3^1+1}
\cdot 7^{2^1\cdot 3^{11} + 1} \cdot 11^{2^3\cdot3^1+1} \cdot
13^{2^1\cdot3^{10}+1} = \\
2^{13\,123}\cdot 3^{39\,367}\cdot 5^{13}\cdot 7^{354\,295}\cdot11^{25}\cdot13^{118\,099}.
\end{multline*}
\end{ex}

\end{document}
